\section{Cyclic layers: a signed comparison beyond the bipartite case}
\label{sec:periodiccomparison}
The bipartite cancellation is one case of a finite cyclic-layer
identity. Here the layers are part of the observable; they are
not discarded by replacing the kernel with a norm envelope.

\begin{theorem}[Relative comparison for a fixed number of cyclic layers]
\label{thm:periodiccomparison}
Fix an integer \(k\geq2\) and \(\delta\in(0,1)\).
For \(i\) modulo \(k\), let \(T_i\) be real nonnegative doubly
stochastic matrices of order \(m\), with
\[
 B_i=T_i-J_m/m,\qquad \|B_i\|\leq1-\delta,\qquad
 \max T_i\leq A/m,\qquad 2\leq A=o(\log m).
\]
Let \(P\) have block \(T_i\) from layer \(i\) to layer \(i+1\)
and all other blocks zero. Put \(C=B_1B_2\cdots B_k\).
Then
\[
 \operatorname{cyc}P
 =\frac{(m!)^k}{m\,m^{km}}\,
   \frac{\det(I-C)}
        {\prod_{i=1}^k\sqrt{\det(I-B_i^\top B_i)}}\,
   \left(1+O_{\delta,k}
       \left(m^{-1/9}\exp(C_{\delta,k}A)\right)\right).
\]
The positive determinant factors are retained in full.
No commutativity, common eigenbasis, or matching row and
column supports is assumed.
\end{theorem}
\begin{proof}
In Liu's exact identity a selected principal submatrix has
zero determinant unless its selected cardinality is the
same, say \(r\), in every layer.
For selected sets \(I_1,\ldots,I_k\), each of size \(r\),
the determinant is exactly
\[
 (-1)^{(k-1)r}\prod_{i=1}^k
                    \det T_i[I_i,I_{i+1}],\qquad I_{k+1}=I_1.
\]
The complementary permanent is the product of the \(k\)
corresponding complementary permanents.
Combining the sign with \((-1)^{kr-1}\) and the coefficient
\(kr\), and dividing by \(k\), gives
\[
 m\operatorname{cyc}P
 =\sum_{r=1}^m(-1)^{r-1}r
  \sum_{\substack{I_1,\ldots,I_k\subseteq[1,m]\\|I_i|=r}}
  \prod_{i=1}^k
   \left(\det T_i[I_i,I_{i+1}]\,
        \operatorname{per}T_i[I_i^c,I_{i+1}^c]\right).
 \tag{20}
\]
This identity includes every compensation term on mismatched
supports.

Normalize by \(\prod_i\operatorname{per}T_i\).
For \(r\leq(\log(2m))^3\), each complementary-permanent
ratio is \(1+O_\delta(m^{-1/9})\), by
Lemma~\ref{lem:relativecompensatedminors}.
Their product has the same error with a constant depending on \(k\).
To bound only that error, let \(M_i\) be the matrix indexed
by \(r\)-subsets, with entries
\(|\det T_i[I,J]|\). Cauchy--Binet gives
\[
 \|M_i\|_F^2=e_r(T_iT_i^\top)\leq A^r/r!.
\]
Repeated Hilbert--Schmidt inequalities, for \(k\geq2\), imply
\[
 \sum_{I_1,\ldots,I_k}\prod_i|\det T_i[I_i,I_{i+1}]|
   =\operatorname{tr}(M_1\cdots M_k)
   \leq\prod_i\|M_i\|_F
   \leq(A^r/r!)^{k/2}.
\]
In particular the total error mass is at most
\[
 \sum_{r\geq1}r(A^r/r!)^{k/2}
               \leq(Ae^A)^{k/2}\leq\exp(C_kA).
\]
For the middle inequality, use \(k/2\geq1\) and
\(\sum_r rA^r/r!=Ae^A\).
Absolute values have been used only after separating the
minor-ratio error.

For the retained tail \(r>(\log(2m))^3\), the full product
of permanent ratios is at most \(\exp(kr+C_kA\log(2m))\).
The same Hilbert--Schmidt bound and \(r!\geq(r/e)^r\)
therefore give a total tail at most
\[
 e^{C_kA\log(2m)}
 \sum_{r>(\log(2m))^3}r
          (e^3A/r)^{kr/2}
 \leq \exp(-c_k(\log(2m))^3\log\log(2m)).
\]
The factor \(e^3\) retains the full \(e^{kr}\) permanent-ratio
cost. The determinant-only tail is smaller.

The retained signed main term is evaluated by the
Cauchy--Binet formula for products:
\[
 \sum_{I_1,\ldots,I_k}\prod_i\det T_i[I_i,I_{i+1}]
                       =e_r(T_1\cdots T_k).
\]
Since \(T_1\cdots T_k=J_m/m+C\),
\[
 \det(I-xT_1\cdots T_k)=(1-x)\det(I-xC).
\]
Differentiating at \(x=1\) now gives
\[
 \frac{m\operatorname{cyc}P}{\prod_i\operatorname{per}T_i}
 =\det(I-C)+
 O_{\delta,k}\left(m^{-1/9}e^{C_kA}
       +e^{-c_k(\log(2m))^3\log\log(2m)}\right).
\]
This is the signed step; bounding the main term by the
Hilbert--Schmidt mass would not prove the formula.

Finally put \(\rho=1-\delta\).
For \(j\geq1\), the two Frobenius factors and the remaining
operator factors give
\[
 |\operatorname{tr}(C^j)|\leq A\rho^{kj-2}.
\]
The logarithmic determinant series is absolutely convergent,
so \(|\log\det(I-C)|\leq C_{\delta,k}A\).
Its determinant is positive by continuity from \(I\) and
\(\|C\|<1\). Dividing by it costs at most \(e^{C_{\delta,k}A}\).
Apply the growing-entry permanent theorem to every \(T_i\).
All full determinants remain in the displayed result,
and all relative errors fit the stated bound.
\end{proof}

\begin{corollary}[Why two layers are special]
\label{cor:twolayersspecial}
If \(k=2\) and \(T_2=T_1^\top\), the spectral factor in
Theorem~\ref{thm:periodiccomparison} is exactly one.
For \(k\geq3\), if all cross blocks equal a symmetric matrix
\(T=J_m/m+B\), that factor is
\[
 \frac{\det(I-B^k)}{\det(I-B^2)^{k/2}}.
\]
Thus higher layers retain a genuine joint spectral contribution.
\end{corollary}
\begin{proof}
For two layers \(C=B_1B_1^\top\); the numerator cancels
the product of the two square-root determinants.
For equal symmetric blocks \(C=B^k\), while each denominator
is \(\sqrt{\det(I-B^2)}\).
These are exact finite identities, preceding any limiting estimate.
\end{proof}

\begin{remark}
The new comparison range is \(A=o(\log m)\), with fixed \(k\).
There is no claim of uniformity when the number of layers grows.
The proof neither invokes an aperiodic hypothesis on the full
\(km\)-vertex matrix nor removes its \(k\) exact unit-circle modes.
They are encoded by the cyclic block structure from the outset.
\end{remark}
